% CONTRACT
%   Establishes: why validation checks in this class of workflow can pass while
%   the property they name goes untested. Introduces the decomposition only as
%   far as the audit needs it.
%   Sources: no result tables — this section carries no measured numbers.
%   May not claim: anything about silencing, and any field-wide prevalence rate.
%     The 21 ledger entries are a starting taxonomy, NOT a measured field-wide
%     rate (decision record §5).
%   Must contain: an explicit statement that the adenoma example is subordinate.

\subsection{The estimand, and the third verdict}
A change in a differentiation marker between two conditions can be written as the
sum of a compositional term (the fraction of a cell type changed), an intrinsic
term (the per-cell output changed), and an interaction
\cite{kitagawa1955,blinder1973,oaxaca1973}. The decomposition is useful only if
each term can be estimated from the data at hand, and for many samples it cannot:
too few cells of the relevant type survive to ask. We therefore report three
outcomes rather than two, and an unestimable intrinsic term is recorded as not
estimable rather than as zero. Conflating the two turns an absent measurement
into a measurement of absence, and every aggregate that absorbs it inherits the
error.

\subsection{The problem: a check that cannot fail}
A validation check states a failure condition. Whether it can detect that
condition is a separate question from whether it passes, and a passing check
answers only the second. The failure mode is silent: an interval reported at 95\%
that covers less often; a calibration reported at a threshold the grid cannot
reach; a guard whose fixture supplies the quantity it is supposed to test. None
raises. This is the shape we are concerned with, and it is not specific to this
workflow.

\subsection{What this paper does}
We audit one differential-composition workflow by executable counterexample: for
each check, construct the input that must make it fail, run it, and measure the
consequence. Where a check cannot be made to fail, that is the finding. We report
the audit's measured consequences, including a failure the audit found in its own
infrastructure, and we apply the same procedure to the workflow's descriptive
result to establish which of its estimates remain reportable. The colorectal
observation is the worked example that makes the audit concrete; it is
subordinate to the validation question and is not the contribution.

\subsection{Related work}
Simulation studies are routinely reported against frameworks such as ADEMP, which
ask that aims, data-generating mechanism, estimands, methods and performance
measures be stated in advance \cite{morris2019}; plasmode simulation draws the
data-generating mechanism from real data \cite{franklin2014}; and congeniality
asks whether the simulation and the analysis are compatible \cite{meng1994}. Our
concern is adjacent but distinct: not whether a simulation is well designed, but
whether a validation check can detect the failure it names. Selective
classification and risk control \cite{geifman2017,elyaniv2010} share the
abstention framing we use for the third verdict.
